%%
%% Blank Springer Nature (sn-jnl) article skeleton
%% Stripped of documentation/instructional text — structure only.
%%
%% Submit as a single .tex file; do not \input other .tex files, and
%% attach figures separately rather than embedding them in the source.
%%
%% Pick ONE \documentclass line matching your target reference style
%% (uncomment it, keep the rest commented/removed):
%%\documentclass[referee,sn-basic]{sn-jnl}          %% "referee" = double line spacing for review
%%\documentclass[lineno,pdflatex,sn-basic]{sn-jnl}  %% margin line numbers
%%\documentclass[pdflatex,sn-nature]{sn-jnl}        %% Nature Portfolio journals
%%\documentclass[pdflatex,sn-basic]{sn-jnl}         %% Basic Springer Nature / Chemistry style
\documentclass[pdflatex,sn-mathphys-num]{sn-jnl}    %% Math & Physical Sciences, numbered refs
%%\documentclass[pdflatex,sn-mathphys-ay]{sn-jnl}   %% Math & Physical Sciences, author-year refs
%%\documentclass[pdflatex,sn-aps]{sn-jnl}           %% American Physical Society style
%%\documentclass[pdflatex,sn-vancouver-num]{sn-jnl} %% Vancouver, numbered refs
%%\documentclass[pdflatex,sn-vancouver-ay]{sn-jnl}  %% Vancouver, author-year refs
%%\documentclass[pdflatex,sn-apa]{sn-jnl}           %% APA style
%%\documentclass[pdflatex,sn-chicago]{sn-jnl}       %% Chicago humanities style
%% Note: sn-basic.bst / sn-chicago.bst support adding "Numbered" in the
%% class options parenthesis to switch between name-date and numbered citation styles.

\usepackage{graphicx}%
\usepackage{multirow}%
\usepackage{amsmath,amssymb,amsfonts}%
\usepackage{amsthm}%
\usepackage{mathrsfs}%
\usepackage[title]{appendix}%
\usepackage{xcolor}%
\usepackage{textcomp}%
\usepackage{manyfoot}%
\usepackage{booktabs}%
\usepackage{algorithm}%
\usepackage{algorithmicx}%
\usepackage{algpseudocode}%
\usepackage{listings}%

%% ---------------------------------------------------------------
%% Theorem-like environments. Three predefined styles:
%% thmstyleone   - numbered, bold head, italic body
%% thmstyletwo   - numbered, roman head, italic body
%% thmstylethree - numbered, bold head, roman body
%% ---------------------------------------------------------------
\theoremstyle{thmstyleone}%
\newtheorem{theorem}{Theorem}%              %% continuous numbering
%%\newtheorem{theorem}{Theorem}[section]%   %% sectionwise numbering instead
\newtheorem{proposition}[theorem]{Proposition}%  %% shares numbering with theorem
%%\newtheorem{proposition}{Proposition}%    %% use for independent numbering instead

\theoremstyle{thmstyletwo}%
\newtheorem{example}{Example}%
\newtheorem{remark}{Remark}%

\theoremstyle{thmstylethree}%
\newtheorem{definition}{Definition}%

\raggedbottom
%%\unnumbered%  %% uncomment for unnumbered section heads

\begin{document}

\title[]{}

%% ---------------------------------------------------------------
%% Author name parts:
%%   GivenName  -> \fnm{}
%%   Particle   -> \spfx{}   (surname prefix, e.g. "van der")
%%   FamilyName -> \sur{}
%%   Suffix     -> \sfx{}    (e.g. "IV")
%% \author*[...] marks the corresponding author; the numeric list in [ ]
%% ties each author to one or more \affil[...] blocks below.
%% Use \equalcont{} for an "these authors contributed equally" note.
%% ---------------------------------------------------------------
\author*[1,2]{\fnm{} \sur{}}\email{}

\author[2,3]{\fnm{} \sur{}}\email{}
%%\equalcont{These authors contributed equally to this work.}

%% Repeat \author[...]{...}\email{...} for each additional author.

\affil*[1]{\orgdiv{}, \orgname{}, \orgaddress{\street{}, \city{}, \postcode{}, \state{}, \country{}}}

\affil[2]{\orgdiv{}, \orgname{}, \orgaddress{\street{}, \city{}, \postcode{}, \state{}, \country{}}}

%% ---------------------------------------------------------------
%% Abstract — choose ONE form and delete the other. Unstructured is
%% the default; structured (with subheadings) is only for journals
%% that explicitly permit it — check the target journal's guidelines
%% for word limits and whether subheadings/equations/citations are allowed.
%% ---------------------------------------------------------------
\abstract{}

%%\abstract{\textbf{Purpose:} \\
%%\textbf{Methods:} \\
%%\textbf{Results:} \\
%%\textbf{Conclusion:} }

\keywords{}

%% Optional classification codes for certain disciplines:
%%\pacs[JEL Classification]{}
%%\pacs[MSC Classification]{}

\maketitle

\section{Introduction}\label{sec:intro}
%% No subheadings in the Introduction; some overlap with the abstract is fine.

\section{Results}\label{sec:results}

\section{}\label{sec:generic}

\subsection{}\label{subsec:generic}

\subsubsection{}\label{subsubsec:generic}

\section{Equations}\label{sec:equations}

%% Inline math: $ ... $
%% Numbered display equation:
\begin{equation}
\label{eq:1}
\end{equation}

%% Multi-line aligned equations — use \nonumber on every line except the
%% last one that should carry a number; \label only on that last line.
\begin{align}
 &  \nonumber \\
 & \label{eq:2}
\end{align}

\section{Tables}\label{sec:tables}

%% Standard table with booktabs rules. Use \footnotetext[]{} for table
%% footnotes (rendered just below the table) and \footnotemark[] to mark
%% the reference point in the table body.
\begin{table}[h]
\caption{}\label{tab:1}%
\begin{tabular}{@{}llll@{}}
\toprule
 & & & \\
\midrule
 & & & \\
\botrule
\end{tabular}
%%\footnotetext{}
\end{table}

%% Full-textwidth table (tabular*) — use when a table is too wide for a
%% single column:
%%\begin{table}[h]
%%\caption{}\label{tab:2}
%%\begin{tabular*}{\textwidth}{@{\extracolsep\fill}lcccccc}
%%\toprule
%%\midrule
%%\botrule
%%\end{tabular*}
%%\end{table}

%% Double-column layout: use table* instead of table for tables spanning
%% both columns. For tables too long to fit in textwidth, use the
%% sidewaystable (single column) or sidewaystable* (double column)
%% environment instead of table*.
%%\begin{sidewaystable}
%%\caption{}\label{tab:3}
%%\begin{tabular*}{\textheight}{@{\extracolsep\fill}lcccccc}
%%\toprule
%%\midrule
%%\botrule
%%\end{tabular*}
%%\end{sidewaystable}

\section{Figures}\label{sec:figures}
%% Use .eps for LaTeX compilation, .pdf/.jpg/.png for PDFLaTeX. Each
%% figure should come from a single image file; avoid subfigures.
\begin{figure}[h]
\centering
\includegraphics[width=0.9\textwidth]{}
\caption{}\label{fig:1}
\end{figure}
%% For a double-column spanning figure, use figure* instead of figure.

\section{Algorithms, Program Code, and Listings}\label{sec:algorithms}

%% Algorithm environment (algorithm/algorithmicx/algpseudocode packages):
\begin{algorithm}
\caption{}\label{algo:1}
\begin{algorithmic}[1]
\Require
\Ensure
\State
\end{algorithmic}
\end{algorithm}

%% Program code via verbatim:
%%\begin{verbatim}
%%
%%\end{verbatim}

%% Code listings via the listings package (\lstset{} configures language,
%% style, frame, etc. per block):
%%\lstset{}
%%\begin{lstlisting}
%%
%%\end{lstlisting}

\section{Cross Referencing}\label{sec:crossref}
%% \label{} goes inside/just below \caption{} for figures/tables, and on
%% the last (non-\nonumber) line of an align block for equations. Use
%% \ref{} to cross-reference; e.g. Figure~\ref{fig:1}, Table~\ref{tab:1},
%% line~\ref{algln} of Algorithm~\ref{algo:1}.

\subsection{Reference Citations}\label{subsec:citations}
%% This template uses natbib to support both numeric and author-year
%% citation styles (selected via the documentclass option above).
%% \cite{} and \citep{} behave differently depending on style — see the
%% template's User Manual for details.

\section{Theorem-like Environments}\label{sec:theorems}

\begin{theorem}[]\label{thm:1}

\end{theorem}

\begin{proposition}

\end{proposition}

\begin{example}

\end{example}

\begin{remark}

\end{remark}

\begin{definition}[]

\end{definition}

%% Predefined proof environment — prints "Proof" in italic, body in
%% roman, with a Q.E.D. box at the end:
\begin{proof}

\end{proof}

%% To label a proof by the theorem it proves:
%%\begin{proof}[Proof of Theorem~{\upshape\ref{thm:1}}]
%%
%%\end{proof}

%% Quote environment:
%%\begin{quote}
%%
%%\end{quote}

\section{Methods}\label{sec:methods}
%% Topical subheadings are allowed here. Include enough experimental/
%% characterization detail for reproducibility.

%% Required (where applicable) for work on live vertebrates, humans, or
%% human samples — an explicit statement covering:
%%   1. Approval (named institutional/licensing committee)
%%   2. Accordance (methods followed relevant guidelines/regulations)
%%   3. Informed consent (for human subjects/tissue)
%%\textbf{Ethical approval declarations}

\section{Discussion}\label{sec:discussion}
%% "Discussion" and "Conclusion" are often interchangeable across
%% disciplines; not mandatory to use both — check journal guidance.

\section{Conclusion}\label{sec:conclusion}

\backmatter

\bmhead{Supplementary information}
%% State here if the article has accompanying supplementary files.

\bmhead{Acknowledgements}
%% Not compulsory; keep brief if included. Grant/contribution numbers
%% may be acknowledged here.

%% ---------------------------------------------------------------
%% Declarations — some journals require this standardized section.
%% If a subsection is not relevant, keep the heading and write
%% "Not applicable" rather than omitting it.
%% ---------------------------------------------------------------
\section*{Declarations}

\begin{itemize}
\item Funding
\item Conflict of interest/Competing interests
\item Ethics approval and consent to participate
\item Consent for publication
\item Data availability
\item Materials availability
\item Code availability
\item Author contribution
\end{itemize}

%% ---------------------------------------------------------------
%% Appendices — use the "appendices" environment (from the appendix
%% package), not the bare \appendix command.
%% Nature Portfolio submissions should title this "Extended Data" instead.
%% ---------------------------------------------------------------
\begin{appendices}

\section{}\label{app:1}

\end{appendices}

%% ---------------------------------------------------------------
%% Bibliography. For Nature Portfolio eJP submissions, paste the
%% .bbl content directly into this file and remove \bibliography{}.
%% ---------------------------------------------------------------
\bibliography{sn-bibliography}
%%\input sn-article.bbl

\end{document}
%% End of blank Springer Nature (sn-jnl) skeleton.
